\subsection{Unary Operations}
\label{subsec:unary-operations}

\begin{tcolorbox}[title={Definition --- Unary Operation}]
\begin{definition}[Unary Operation]
\label{def:unary-operation}
Let $S$ be a set.
A \textbf{unary operation} on $S$ is a function
\[
f:S \to S.
\]
\end{definition}
\end{tcolorbox}

\begin{remark*}[Standard quantified statement]
\[
\operatorname{UnaryOperation}(f,S)
\Longleftrightarrow
f:S\to S.
\]
\end{remark*}

\begin{remark*}[Definition predicate reading]
$f$ is a unary operation on $S$ exactly when it takes one element of $S$ and
returns one element of $S$.
\end{remark*}

\begin{remark*}[Negated quantified statement]
\[
\neg\operatorname{UnaryOperation}(f,S)
\Longleftrightarrow
\exists x\in S,\; f(x)\notin S.
\]
\end{remark*}

\begin{remark*}[Interpretation]
Unary operations appear throughout algebra: negation $a \mapsto -a$,
multiplicative inverse $a \mapsto a^{-1}$, complex conjugation
$z \mapsto \bar z$, and set complementation are all one-input operations.
\end{remark*}

\begin{remark*}[Dependencies]
\begin{itemize}
  \item Function.
  \item Carrier set.
\end{itemize}
\end{remark*}
